% ---
% id: meat-core/lem:product-marginal
% kind: lem
% title: "Product marginal of Alice's retained systems"
% summary: >-
%   In trusted-source mode, under retention, emission of the specified
%   ensemble, sequential rounds, private source randomness and a sealed
%   laboratory, the n rounds start with Alice's systems in the declared
%   marginal in product, and each round finds her system in that marginal.
% prominence: standard
% status: accepted
% uses: [def:roundmodel, transmitter/def:source-spec, transmitter/def:srcert, transmitter/def:source-segment, transmitter/thm:certificate-sufficiency]
% mentions: [thm:accumulated-entropy, prop:trivial-alice-system, transmitter/ass:src, transmitter/ass:ens, transmitter/ass:seq, transmitter/decoy/ass:wcp]
% assumes: [A-SRC, A-ENS, A-SEQ, A-SEED, A-LAB]
% refs: []
% provenance: newly-derived
% lean:
% notes: "Parts (ii) and (iii) state the product marginal of A (omega_{A_1^n} = tau_A^{otimes n}, per round omega_A = tau_A, in trusted-source mode), derived here under A-SRC, A-ENS, A-SEQ, A-SEED (v) and A-LAB (i) and (ii) at the party of A, with the adversary's initial register finite-dimensional, as every register of the round model is (a hypothesis of the statement, not a premise about the adversary: an adversary with memory of any dimension is reached by cor:composed-security through lem:memory-lift). The untrusted-mode reading is stated where it is used (thm:accumulated-entropy, fact:meat locator note (vii)), and the d_A = 1 case is the last sentence of the statement."
% ---

In \tref{trusted-source mode} the \tref{round model} takes Alice's retained
systems as inputs of the round channels, and the accumulation theorem
(Theorem~\ref{thm:accumulated-entropy}) asks that they start in the declared
marginal, in product over the rounds. The lemma derives that from the
transmitter's premises. The \tref{source-replacement certificate} describes
one round's preparation as Alice's measurement of her half of the certificate
(\xref{transmitter/thm:certificate-sufficiency}\,(iii)); A-ENS
(\xref{transmitter/ass:ens}) makes the $n$ preparations a product, and
A-SEED (Assumption~\ref{ass:seed}\,(v)) makes them independent of what the
adversary holds at the start; and retention, A-SRC
(\xref{transmitter/ass:src}), the sealed laboratory of the party that holds
$A$, A-LAB (Assumption~\ref{ass:lab}\,(i), (ii)), and the sequential round structure,
A-SEQ (\xref{transmitter/ass:seq}), let each measurement be moved into its
round. Retention alone does not give the product: it fixes what may
act on each round's system, and the independence of the source's rounds gives
the product. A specification declared by a \tref{source model} carries, in
the place of A-ENS, the premise its submodule declares, which refines A-ENS:
for weak coherent pulses it is A-WCP (\xref{transmitter/decoy/ass:wcp}), whose
last three clauses are A-ENS's for the model's states. At $d_A=1$ the conclusion holds of every state
(Proposition~\ref{prop:trivial-alice-system}\,(ii)).

\begin{lemma}\label{lem:product-marginal}
Let the declaration be in \tref{trusted-source mode} with $d_A\ge2$, and let
its $A$, $\tau_A$ and $\{M_x\}$ be those of a \tref{source specification}
(\xref{transmitter/def:source-spec}) that carries a source-replacement
certificate $\Psi$ for $(A,\mathbb{T},\tau_A,\{M_x\})$ against its
\tref{virtual ensemble} $\{p_x,\rvirt_x\}$ on $\Avirt$, in any of the forms of
\xref{transmitter/def:srcert}, an \tref{ensemble presentation} being read as
the block-form certificate it defines; let its \tref{source segment} satisfy
the source identity $\Xi(\rvirt_x)=\rho_x$ for every $x\in\Xset$, with $\Xi$
the composite channel of the segment and $\Xi=\idmap$ when the segment is
empty (\xref{transmitter/def:source-segment}\,(S3)). Let A-SRC
(\xref{transmitter/ass:src}), A-ENS (\xref{transmitter/ass:ens}), A-SEQ
(\xref{transmitter/ass:seq}), A-SEED (Assumption~\ref{ass:seed})
and A-LAB (Assumption~\ref{ass:lab}) at the party to which the
\tref{party partition} assigns $A$ hold of the physical protocol, let the
adversary's initial register $E$ be finite-dimensional, and let $\sigma_E$
be the state in which it starts. Put
\begin{equation}\label{eq:product-input}
  \omega^{0}\coloneqq\Big(\bigotimes_{j=1}^{n}\Psi_{A_j\Avirt_j}\Big)
  \otimes\sigma_E\ \in\ \Seq{A_1^nE_0},
  \qquad E_0\coloneqq(\Avirt)_1^n\,E .
\end{equation}
Then:
\begin{enumerate}[label=(\roman*),leftmargin=2.2em]
\item \emph{(source replacement over the $n$ rounds)} against every attack,
the physical protocol and the protocol that starts from $\omega^{0}$, applies
$\Xi$ to $\Avirt_j$ in round $j$ in the place of the source's emission, and
applies $\pcNmsm$, Alice's measurement as the channel of
\xref{transmitter/thm:certificate-sufficiency}\,(iii), to $A_j$ at the point of
round $j$ at which the declared protocol's measurements are made, produce
the same state after every round, on every register that both then hold; in
the second protocol $A_j$ is an input of that measurement and of no other
operation, which is how the round model of Definition~\ref{def:roundmodel}
takes Alice's system in trusted-source mode;
\item $\omega^{0}_{A_1^n}=\tau_A^{\otimes n}$;
\item for every \tref{attack sequence} of Definition~\ref{def:roundmodel}
and every $j$, the state on which the round channel $\Mprot_j$ acts, in a run
of the round model from $\omega^{0}$, has the marginal $\tau_A$ on $A_j$.
\end{enumerate}
At $d_A=1$, where $\tau_A=[1]$ and no certificate is declared, (ii) and (iii)
hold of every state on $A_1^nE_0$.
\end{lemma}

\begin{proof}
\emph{Idea.} The physical preparation is a product over the rounds of one
round's classical-quantum source state, which the certificate writes as
Alice's measurement of her half of $\Psi$ followed by $\Xi$; nothing touches
her half or her letter before the round's measurements, so the interchange
law moves her measurement there; and the marginal of a product of copies of
$\Psi$ on Alice's side is the product of their marginals.

\emph{At $d_A=1$.} $\Seq{A}=\{[1]\}$, so $\Seq{A_1^n}$ is a singleton and
every state on $A_1^nE_0$, and every state the round channels act on, has the
marginal $[1]^{\otimes n}=\tau_A^{\otimes n}$, respectively $[1]=\tau_A$.

\emph{Step 1 (the physical source state).} Let $X_j$ hold the letter of
round $j$, and write
$\rho_{X\Asig}\coloneqq\sum_{x\in\Xset}p_x\dyad{x}_X\otimes\rho_x$ for one
round's classical-quantum source state on the \tref{transmitted system}. By
A-ENS (i) and (ii) the letters of the $n$ rounds are independent with the
weights $p_x$ and, given them, the emitted systems are in
$\rho_{x_1}\otimes\dots\otimes\rho_{x_n}$, so the letters and the emitted
systems are jointly in
$\sum_{x_1^n}\prod_jp_{x_j}\dyad{x_1^n}\otimes\rho_{x_1}\otimes\dots\otimes
\rho_{x_n}=\rho_{X\Asig}^{\otimes n}$. By A-ENS (iii) neither the letter of
round $j$ nor the state emitted in it depends on anything that happens in the
other rounds, so preparing all $n$ letters and systems at the start in that
state, and releasing the round-$j$ system only in round $j$, changes nothing
the protocol produces. By A-SEED (v) the letters and the prepared systems
are jointly independent of the adversary's initial register, and by its last
clause no purification of a prepared system is made available to the
adversary during or after the protocol, so that no register beyond the
letters and the emitted systems carries the preparation; the protocol
starts from $\rho_{X\Asig}^{\otimes n}\otimes\sigma_E$.

\emph{Step 2 (the certificate).} By the source identity,
$\rho_{X\Asig}=(\idmap_X\otimes\Xi)(\rho_{X\Avirt})$ with
$\rho_{X\Avirt}\coloneqq\sum_xp_x\dyad{x}_X\otimes\rvirt_x$. By (SR2), in the
form of \xref{transmitter/thm:certificate-sufficiency}\,(iii)(b) with the
identity channel in the place of $\mathcal{E}$,
$(\pcNmsm\otimes\idmap_{\Avirt})(\Psi)=\rho_{X\Avirt}$. Hence
$\rho_{X\Asig}=(\pcNmsm\otimes\Xi)(\Psi)$, and the initial state of Step~1 is
$\big(\bigotimes_j(\pcNmsm\otimes\Xi)_{A_j\Avirt_j}\big)(\omega^{0})$: the
physical protocol is the protocol that starts from $\omega^{0}$ and applies,
at the start, $\pcNmsm$ to every $A_j$ and $\Xi$ to every $\Avirt_j$, the
round-$j$ output of $\Xi$ being released in round $j$. Applying $\Xi$ to
$\Avirt_j$ at the start or in round $j$ produces the same state, since
nothing else acts on $\Avirt_j$ before round $j$.

\emph{Step 3 (moving the measurement into the round).} Fix $j$. Let
$\mathcal{E}_j$ be the composite of every operation of the protocol, of the
parties and of the adversary, from the start to the point of round $j$ at
which the declared protocol's measurements are made. It acts on neither $A_j$
nor $X_j$; the measurements of the other rounds, wherever they stand, act on
$A_k$ and $X_k$ with $k\ne j$. On $A_j$: by A-SRC (ii) no operation of the
protocol or of the adversary acts on $A$ except the declared measurement on it.
On $X_j$, for the adversary's operations: the letters are records of the
party to which the party partition assigns $A$, and by A-LAB (i) at that party
no information about them leaves its laboratory except through the systems
handed out of it, which in trusted-source mode are the systems its source
emits, and the public messages of the protocol, and by A-LAB (ii) the
adversary acts on its laboratory only through the systems delivered into it,
as its declared measurements receive them, and the messages the protocol
delivers to it, so that nothing she does reads or writes the letter record;
no operation of the adversary acts on $X_j$
itself: what the emitted systems carry about the
letter is already in the state of Step~2, and the messages are treated next. On $X_j$,
for the protocol's own operations: the letter of round $j$ is a record of
that same party, which the declared protocol
reads only in the round-$j$ protocol map, at or after the round's
measurements, and in the post-processing; by A-SEQ (i) the messages of round
$j$ are sent after every measurement of the round, by A-SEQ (iv) no operation
on a quantum system of round $j$ depends on them, and by A-SEQ (ii) the
post-processing follows all $n$ rounds. So $\pcNmsm$, a map out of $A_j$
alone into $X_j$, and $\mathcal{E}_j$ act on complementary registers, and by
the interchange law of \xref{transmitter/thm:certificate-sufficiency}\,(iii)(a)
$\mathcal{E}_j\circ(\pcNmsm)_{A_j}=(\pcNmsm)_{A_j}\circ\mathcal{E}_j$ as maps on
the whole input. Applying this for $j=1,\dots,n$ moves each measurement from
the start to its round without changing any state the protocol produces after
the round, which is~(i); in the protocol so obtained $A_j$ enters no operation
but the round-$j$ measurement.

\emph{Step 4 (the marginals).} (ii): $\sigma_E$ is a state, so
$\omega^{0}_{A_1^n}=\bigotimes_j\tr_{\Avirt_j}\Psi_{A_j\Avirt_j}
=\tau_A^{\otimes n}$ by (SR1). (iii): in a run of the round model from
$\omega^{0}$, the state on which $\Mprot_j$ acts is produced from $\omega^{0}$
by $\Mprot_1,\dots,\Mprot_{j-1}$ and the attack channel $\Aev_j$
(Definition~\ref{def:roundmodel}, \eqref{eq:roundchannel}), and in
trusted-source mode each $\Mprot_k$ acts on $A_kE_{k-1}$ and $\Aev_j$ on
$E_{j-1}$, none of them on $A_j$. A channel on a register complementary to
$A_j$ leaves the $A_j$-marginal unchanged, the partial trace over that
register being invariant under it, so the marginal is $\omega^{0}_{A_j}$, which
is $\tau_A$ by~(ii).
\end{proof}
